%=====================================================================
\section{Introduction}\label{sec:intro}
%=====================================================================

For a prime $p$ and $n\ge1$ put
\[
P_{p,n}(q)=\prod_{\substack{k\le pn\\ p\nmid k}}(1-q^k)=\frac{(q;q)_{pn}}{(q^p;q^p)_n}.
\]
P.~Borwein observed, and Andrews \cite{Andrews} made widely known, that the coefficients of such products and of their
small powers appear to have a sign pattern that depends only on the exponent modulo $p$.
% Andrews [J. Symbolic Comput. 20 (1995)] discusses several conjectures of P. Borwein for p=3 and p=5; the names
% First/Second/Third follow Wang--Krattenthaler and Berkovich--Dhar (First: P_{3,n}; Second: P_{3,n}^2; Third: P_{5,n}).
% TODO: numbering inside Andrews's paper itself not verified (full text not accessible).
The three conjectures that concern us are the following. Write $c(m)$ for the coefficient of $q^m$ in the polynomial
under discussion.
\begin{enumerate}
\item (\emph{First Borwein conjecture.}) For $P_{3,n}$: $c(m)\ge0$ if $m\equiv0$ and $c(m)\le0$ if $m\equiv1,2\pmod 3$.
\item (\emph{Cubic Borwein conjecture.}) The same pattern $+\,-\,-$ holds for $P_{3,n}^3$.
\item (\emph{Third Borwein conjecture.}) For $P_{5,n}$: $c(m)\ge0$ if $5\mid m$ and $c(m)\le0$ otherwise.
\end{enumerate}
There is also a second Borwein conjecture, for $P_{3,n}^2$. (The names First, Second and Third follow
\cite{WK,BerkovichDhar}; Andrews \cite{Andrews} discusses Borwein's conjectures for $p=3$ and $p=5$.) The first conjecture was
proved by Wang \cite{Wang}, by an analytic (saddle-point) method. Wang and Krattenthaler \cite{WK} gave a new proof of the first
conjecture and proved the second by an asymptotic approach to Borwein-type sign patterns; they also proved ``two thirds'' of
the cubic conjecture (see also \cite[\S1]{BerkovichDhar}), namely the class $m\equiv0\pmod3$ and, in the range
$m\le\frac92n^2$ (equivalently, by the palindromic symmetry of $P_{3,n}^3$, the complementary halves), the class
$m\equiv1\pmod3$, and explain \cite[\S11]{WK} why the class $m\equiv2$ resists their method.
% Consistency check (resolves the earlier TODO): P_{3,n}^3 is palindromic of degree 9n^2 = 0 (mod 3), so m -> 9n^2-m
% swaps the classes 1 and 2. Berkovich--Dhar's 'all of kappa_n, half of delta_n, half of theta_n' therefore describes
% class 0 plus one pair of mirror-image halves (class 1 for m <= 9n^2/2 <-> class 2 for m >= 9n^2/2). The half that is
% open must be the 'thin' one: by Jacobi's identity every coefficient of P_infty^3 in class 2 vanishes, so class 2 with
% m <= 9n^2/2 is exactly where the main term cancels -- which is what WK Sect. 11 says resists their method. This is the
% class treated in Theorem A. Theorem number in WK Sect. 10 not checked verbatim (arXiv full text not reachable here).
The Third Borwein conjecture is listed as open in the recent literature \cite[Conj.~1.3]{BerkovichDhar}, where it is noted
that \cite{WK} provides ideas for ``three fifths'' of a proof; we are not aware of a proof in the literature.

\subsection{Main results}

\begin{conjecture}[cubic Borwein conjecture]\label{conj:cubic}
Let $c_n(m)=[q^m]P_{3,n}(q)^3$. For all $n\ge1$ and $0\le m\le\frac92n^2$: $c_n(m)\ge0$ if $m\equiv0$, and $c_n(m)\le0$ if
$m\equiv1,2\pmod 3$.
\end{conjecture}

\begin{theoremA}[the remaining class of the cubic conjecture]
For all $n\ge1$ and all $m\equiv2\pmod3$ with $m\le\frac92n^2$ we have $c_n(m)=0$ if $m<3n$ and $c_n(m)<0$ if $m\ge3n$.
Together with \cite{WK} this proves Conjecture~\ref{conj:cubic}.
\end{theoremA}

\begin{conjecture}[Third Borwein conjecture]\label{conj:third}
Let $c_n(m)=[q^m]P_{5,n}(q)$. For every $n\ge1$ and every $m$, $c_n(m)\ge0$ if $m\equiv0\pmod5$ and $c_n(m)\le0$ if
$m\not\equiv0\pmod 5$.
\end{conjecture}

\begin{theoremB}
Conjecture~\ref{conj:third} holds.
\end{theoremB}

Theorem~A is Theorem~\ref{thm:cubic-main} and Theorem~B is Theorem~\ref{thm:third-main}. Both are computer-assisted.
Theorem~A needs only the one residue class $m\equiv2\pmod3$; Theorem~B is proved for all five classes modulo $5$.

\subsection{Why the remaining cases are hard}
Both polynomials factor as $P_{p,n}^\delta=G_p^\delta E_n$ (\S\ref{sec:framework}), where $G_p=(q;q)_\infty/(q^p;q^p)_\infty$
and $E_n$ is a tail product with non-negative coefficients. The circle method gives the coefficients of $G_p^\delta$ with a
main term from the cusps of denominator $p$, and the tail $E_n$ is a small perturbation of $1$ for exponents $m$ below
$N=pn+1$. The difficulty is that in the classes we must treat, the main term of $G_p^\delta$ vanishes:
\begin{itemize}
\item for $p=3$, $\delta=3$, all coefficients of $G_3^3$ in the class $m\equiv2$ vanish (Lemma~\ref{lem:vanish});
\item for $p=5$, $\delta=1$, all coefficients of $G_5$ in the classes $m\equiv3,4$ vanish (Lemma~\ref{lem:vanish}).
\end{itemize}
In these ``thin'' classes $c_n(m)$ is produced entirely by the tail $E_n$, and the leading saddle-point contribution
cancels. The sign then comes from a secondary term that is exponentially smaller than the naive main term.

\subsection{Overview of the method}
The coefficient range is split into four regions (\S\ref{sec:regions}).
\begin{enumerate}
\item \emph{Exact range.} Small $n$ are handled by exact integer computation.
\item \emph{Band.} For $m$ below a small multiple of $N$ only a few parts of $E_n$ contribute, and $c_n(m)$ is a short,
explicit combination of coefficients of $G_p^\delta$, controlled by exact values and by the circle-method expansion.
\item \emph{Laplace region.} Here an exact Bessel--Laplace (or Poisson) representation of the main term is combined with
the closed form of $\log E_n$ at a twisted point, and a saddle-point analysis is carried out in scaled variables.
\item \emph{Centre.} Near the middle of the coefficient range, Cauchy's formula for the polynomial is evaluated directly
by steepest descent through the saddles near the primitive $p$-th roots of unity.
\end{enumerate}
The method was developed on the cubic case. For $p=5$ three new difficulties arise:
\begin{itemize}
\item \emph{two} vanishing classes instead of one;
\item in the Laplace region for the classes $3,4$ the surviving term carries a factor $e^{-5nu}$ which is of size
$e^{O(N)}$ \emph{in the exponent}; the saddle must be re-centred on this ``thin factor'' (\S\ref{sec:lap34});
\item in the centre, the contribution of the circle away from the four peaks must be bounded by an explicit
\emph{off-peak lemma} for $|P_{5,n}|$ (\S\ref{sec:offpeak}), which is then combined with the near-peak analysis box by box.
\end{itemize}
The two proofs share the method, not their results: the cubic theorem is not an ingredient of the proof for $p=5$, and
conversely.

\subsection{What is computer-assisted, and how it is verified}
Each analytic region is reduced to finitely many inequalities between explicit functions of one to three real parameters
(a scaled saddle point, the reciprocal $1/N$, an angle), and these are verified on coverings of the parameter domains by
balls in Arb ball arithmetic \cite{Arb}. Uniformity in $n$ is obtained in two ways: the scaled integrands are analytic in
$1/N$ at $0$, so one ball covers all large $N$; and every term that depends on $n$ is shown, by an asserted inequality, to
be non-increasing in $n$. Exact ranges are computed with FLINT \cite{FLINT} in $\Z[q]$.

Throughout, a statement proved by a computation is stated as a lemma or a certificate, followed by a line
\emph{Certificate:} naming the program and its log. Section~\ref{sec:computations} lists, for every certificate, the
command, the log and the line the log must contain; Appendix~\ref{app:verification} summarises the independent audits of the
proof of Theorem~B. The trust base consists of the written arguments of this paper, the standard facts of
Appendix~\ref{app:facts}, and the correctness of Arb and FLINT.

\subsection{Organisation}
Section~\ref{sec:framework} sets up the common framework and notation. Section~\ref{sec:cubic} proves Theorem~A and
Section~\ref{sec:third} proves Theorem~B. Section~\ref{sec:computations} describes the certificates.
